Actor 将一个状态持有者与送往它的消息结合起来。
外部调用方通过地址或句柄提交消息，
Actor 根据消息处理自己的状态，并可以返回结果或继续发送消息。
这使状态更新集中在明确的执行单元中，
调用方通过消息协作，而不用直接读写另一进程的内存。
具体实现还要规定消息排队、方法并发、实例失败和状态保存规则。

在 Ray 中，Actor 由远程类建立，
创建时为实例安排一个 worker，后续方法发往对应进程。
方法可以访问同一份实例内存，
因此适合保存已加载图像、连接池、模型、
分区内查询索引或一批任务的局部进度。
句柄定位实例，方法调用返回表示未来结果的对象引用。
方法调用关联一次执行和返回结果，句柄则可以继续提交其他方法。
同一实例中的缓存因而能够跨越多次调用复用~\cite{ray259-actors}。

购买系统可以用 Actor 承载某分区的恢复调度：
读取未决 $C/T$，限制并发查询，整理近期响应，再将结果交回权威。
运行时的待执行方法集合类似投往实例的邮箱，
业务权威中的持久待办则记录即使实例消失也要继续承担的工作。
图~\ref{fig:ch20-actor-mailbox} 中，
每次方法先取得自己的原责任，实例内存帮助执行，
原 $S$ 绑定、执行代次和终态继续约束写入。

\CTwentyFigure{ch20-30}{Actor 句柄和待执行调用定位实例，持久待办保存跨实例的责任}{fig:ch20-actor-mailbox}

\subsection{进程重启先发生，状态恢复另有来源}

Ray Actor 的 \code{max\_restarts} 默认 0。
配置正数或无限重启后，意外退出的 Actor 可以重新建立。
运行时重新执行构造函数形成实例。
构造函数若把计数器设为 0，新实例就从 0 开始。
需要接续已处理位置时，
构造函数读取外部检查点或业务权威，再重建内存集合~\cite{ray259-actor-fault}。

设 Actor 内存计数已到 5，随后它完成一次库存预留，
外部事务保存原 $S$、原输入和预留结果，
Actor 却在回复前退出。
新实例内存从 0 开始，重试方法先查询原 $S$，
取得已保存的 $P/C/T$ 和原预留结果。
库存已在原事务里由 1 减至 0，
重试读取原绑定后直接返回，
新实例的计数只反映它自己这次执行。

\CTwentyFigure{ch20-12}{Actor 构造函数重建内存，持久业务记录恢复已提交事实}{fig:ch20-actor-restart}

检查点也要明确保存范围。
若只保存已经处理到第 100 条，
而第 100 条外部调用的结果未知，
恢复者从第 101 条开始就会遗漏一份责任。
检查点可以保存待办身份集合与已确认进度，
外部效果按原身份查询。
检查点与业务结果分别提交时，
恢复步骤要能识别其中一份已保存、另一份尚未保存的状态，
再从原 $C/T$ 决定继续查询还是读取原结果。

检查点所在存储还要能被重启后的实例取得。
新实例可能位于另一节点，
只有旧机器临时目录里的一份文件，
就无法为这次重建提供输入。
本书的完整业务权威和固定图像提供稳定来源，
Actor 检查点减少重扫成本，
权威中的未决集合则决定哪些责任确实没有结束。

\subsection{方法重试与返回结果的缝隙}

Actor 方法的 \code{max\_task\_retries} 默认 0。
实例不可达或进程失败可能让调用方得到错误，
即使方法在失败前已经成功执行。
配置方法重试后，同一个方法可以执行多次。
它对 Actor 内存、外部数据库和 HTTP 接口造成的变化，
分别由这些位置保存的状态决定~\cite{ray259-actor-fault}。

方法若先在内存标记已发送，再调用平台，
重启会丢失标记。
若先调用平台，再保存标记到检查点，
平台成功与检查点写入之间又有一个时间窗口。
支付工作者因此从原 $C/T$ 取得 \code{scope/pay-C}，
发送、查询和回调都汇合到同一权威记录。
重复执行方法时，它先恢复原责任，再取得原结果或接续原操作。

\CTwentyFigure{ch20-13}{方法可能已执行却未返回，重试先恢复原业务身份}{fig:ch20-actor-reply-gap}

普通非 detached Actor 还有 owner，
即最初创建该 Actor 的 worker。
owner 退出后，Actor 与它一同失效，
即使配置了非零 \code{max\_restarts}，也不会自动恢复这份关系。
detached Actor 脱离创建者的生命周期，
可以在原创建者退出后继续存活和按预算重启。
集群销毁、显式销毁与重启预算耗尽，
则会结束相应 Actor 的继续执行条件。
命名 Actor 让其他进程重新取得句柄，
持久业务记录让取得句柄的进程知道应继续哪份责任。

重建可能需要较长时间。
方法在实例尚不可用时继续尝试，会消耗方法重试预算。
恢复调度将健康探测、退避、排队和持久责任扫描结合起来：
先减少发往重建实例的新调用，
探测实例已可用后再逐批取得待办，
便能让构造函数加载图像与后续恢复调用获得各自的资源。

\subsection{顺序由具体 Actor 模型决定}

同步单线程 Actor 在同一提交者、通常执行且没有启用相应乱序选项时，
按该提交者的提交顺序运行方法。
不同提交者之间没有这份共同顺序。
若一个提交者的方法正等待输入依赖，
另一个提交者的就绪方法可能先运行。
方法发生重试或显式允许乱序时，顺序条件也随之改变。
本书按所锁版本的任务顺序说明采用这些条件~\cite{ray259-actor-order}。

异步 Actor 在单线程事件循环中运行协程。
协程是一段可以暂停、再从暂停位置继续的执行过程。
方法在 \code{await} 处等待尚未完成的 I/O 时，
把执行机会交回事件循环，
其他方法便可在同一线程中进入。
这时没有两个 Python 指令在同一时刻占用该线程，
但一次业务判断与它稍后的写入之间已经可以插入另一方法。
线程 Actor 则通过线程池承载多个方法，
共享状态还可能面对线程间的实际并行访问~\cite{ray259-actor-async}。

设一个错误实现只在 Actor 内存保存库存，
初始 \code{available=1}，两个不同购买分别调用方法 A 和 B。
A 先把库存 1 读入局部变量 \code{seenA}，
判断足够，随后 \code{await} 外部查询。
B 在这段等待中进入，也读到 1 并保存 \code{seenB}。
B 查询结束后按 \code{seenB-1} 把库存写成 0，接受购买 B。
A 随后恢复，按它先前保存的 \code{seenA-1} 又写成 0，接受购买 A。
最终内存库存看似为 0，接受记录却已有两份，
一件库存被分配给两笔不同购买。
错误发生在等待前的判断与等待后的写入分离，
这一交错在单线程事件循环中就可以发生。

\CTwentyFigure{ch20-31}{两个协程跨越 await 保存旧观察，库存为 0 却已接受两笔购买}{fig:ch20-actor-await}

如果程序改成恢复后对当前库存直接减一，
两次先前的足量判断也可能先后把库存减为 0 和 $-1$。
两种写法分别暴露丢失更新和超额扣减，
共同缺少的是对接受判断与资源更新的统一仲裁。
Actor 内部可以用适合其执行模型的锁保护完整临界区，
但多实例、重启与持久恢复还要从外部权威取得同一资源事实。

\CTwentyFigure{ch20-14}{同步顺序、协程交错与线程并发分别规定方法怎样进入共享状态}{fig:ch20-actor-modes}

在完整购买模型中，原 $S$ 的输入比较、
库存或额度判断、资源更新与原结果关联，
进入同一权威事务或状态机命令。
以上例为例，库存初始 1，
A 的原子预留先提交后形成可读取的原结果，
B 在自己的仲裁步骤看到库存 0，保存库存不足的终态。
同一 $S$ 重投时读取对应原结果，
Actor 的并发度、实例重建和提交者数量只改变执行安排。

需要顺序的恢复队列还要规定业务依赖。
队列可以按 $C/T$、分区应用位置或明确的前置责任调度。
已完成记录再次到达时取得原终态，
尚未就绪的责任等待自己的依赖。
每个队列项携带自己的 $S$、$C/T$ 与工作代次，
后进入的业务据此取得本项资格。

\subsection{可替换执行者怎样接回完整协议}

Actor 取得恢复工作时，先读取原任务身份与当前执行代次，
按资源端规则申请或核验资格。
调用平台时保留原 \code{scope/pay-C}。
结果回写时，权威核对它属于原 $C/T$，保存完整平台结果。
实例丢失会使尚未确认的工作重新进入恢复集合，
旧实例已经发出的请求则可能继续在网络和平台中运行。

例如，旧 Actor A 的付款请求延迟，
恢复 Actor B 已经开始查询原 $C_1$。
B 暂未查到确定结果时，继续按原键查询或重试允许的原操作。
A 的旧请求随后成功为 $X_1$，
B 再通过原键取得 $X_1$ 并回填。
第十七章的资源端代次核验约束旧写入资格，
平台原身份关联则把已经在途的付款与后继恢复汇合到同一结果。

\CTwentyFigure{ch20-15}{Actor 重建与旧请求在途并存，平台操作沿原 $C$ 汇合}{fig:ch20-actor-inflight}

Actor 关闭时，已经保存在权威中的待办可以由新实例扫描继续。
旧实例停止取得新工作，等待一部分在途调用结束，再退出执行。
只在内存队列中的业务请求则先形成持久受理记录，
再向调用方确认已经接下责任，
新实例才能从保存域重新找到它。
